\clearpage
\subsection{Depth Anything V2 与 Depth Anything V2 (Metric Depth) 的对比分析}

\subsubsection{深度与点云精度比较}

我们在相同的 8 张图上比较两种模型的结果，分别设置 518 和 280 两种输入尺寸。两者采用相同的有效像素、相机内参及对齐、裁剪规则，具体方法见 1.2。同时保留 Depth Anything V2 (Metric Depth) 的原始米制结果，观察其直接测距能力。

\ReportTableCaption{两种模型的深度与点云误差（8 张图等权均值；粗体为同尺寸下两种模型对齐后的最优值）}
\begin{ReportTable}{p{176bp}ccCCC}
\rowcolor{ReportNavy}
\ReportHead{模型} & \ReportHead{尺寸} & \ReportHead{处理} & \ReportHead{AbsRel} & \ReportHead{\shortstack{深度 RMSE\\/ m}} & \ReportHead{\shortstack{XYZ RMSE\\/ m}} \\
Depth Anything V2 & 518 & 对齐后 & \textbf{0.0755} & \textbf{0.3738} & \textbf{0.3986} \\
Depth Anything V2 (Metric Depth) & 518 & 对齐后 & 0.1002 & 0.4792 & 0.5156 \\
Depth Anything V2 (Metric Depth) & 518 & 原始米制 & 0.2222 & 0.7055 & 0.7534 \\
Depth Anything V2 & 280 & 对齐后 & \textbf{0.0804} & \textbf{0.4020} & \textbf{0.4300} \\
Depth Anything V2 (Metric Depth) & 280 & 对齐后 & 0.1190 & 0.6226 & 0.6694 \\
Depth Anything V2 (Metric Depth) & 280 & 原始米制 & 0.5940 & 1.6117 & 1.7377 \\
\end{ReportTable}


\textbf{相同对齐条件下，Depth Anything V2 的三项平均误差在两种输入尺寸下都更低。}说明在本次实验中，能够输出米制深度，并不意味着对齐后的深度和点云一定更准确。

\subsubsection{对齐前后的差距变化}

为直观体现对齐的效果，我们以\textbf{同一输入尺寸下 Depth Anything V2 对齐后的误差为基准}，计算 Depth Anything V2 (Metric Depth) 的误差超出比例，即“（米制深度误差 $\div$ 基准误差 $-1$）$\times100\%$”。以下比例由 8 张图的指标均值计算。

\ReportTableCaption{米制深度误差超出基准的比例（原始米制 $\rightarrow$ 对齐后；粗体为较小差距）}
\begin{ReportTable}{cCCC}
\rowcolor{ReportNavy}
\ReportHead{输入尺寸} & \ReportHead{AbsRel} & \ReportHead{深度 RMSE} & \ReportHead{XYZ RMSE} \\
518 & 194.3\% $\rightarrow$ \textbf{32.8}\% & 88.8\% $\rightarrow$ \textbf{28.2}\% & 89.0\% $\rightarrow$ \textbf{29.4}\% \\
280 & 639.1\% $\rightarrow$ \textbf{48.0}\% & 300.9\% $\rightarrow$ \textbf{54.9}\% & 304.1\% $\rightarrow$ \textbf{55.7}\% \\
\end{ReportTable}


\textbf{对齐后，两种模型的误差差距显著缩小。}在 518 下，深度 RMSE 超出基准的比例从 88.8\% 降至 28.2\%，XYZ RMSE 从 89.0\% 降至 29.4\%；280 下的两项差距也从约 300\% 降至约 55\%。因此，两种模型对齐后的表现更接近，但仍存在一定差距。以米制预测自身为参照，518 下对齐后的深度 RMSE 和 XYZ RMSE 分别降低了 32.1\% 和 31.6\%。

上述变化来自对齐与裁剪的共同作用，原始米制比例仅用于观察处理效果，不作为公平排名。两种模型的训练方式和数据也不同，不能将精度差别全部归因于是否包含尺度信息。

\subsubsection{尺度信息的作用}

结合课上讲的成像关系，在相机内参已知时，\textbf{一个像素只能确定一条从相机光心出发的射线，深度决定三维点位于这条射线的哪一处。}如果缺少深度，射线上多个不同位置的点都可以投影到同一个像素，因此单凭像素位置无法唯一恢复三维点。

尺度信息使重建的深度和空间大小具有实际意义。例如，将场景中的所有三维坐标同时放大两倍，投影到图像上的位置仍可以保持不变，但真实距离和物体大小已经改变。\textbf{尺度信息用于确定真实的距离与大小}，使点云能够用于测距和尺寸估计，而不只是展示相对的前后关系。

在本实验中，Depth Anything V2 的输出没有米制单位，我们利用参考深度进行对齐，为结果补充尺度与偏移信息；Depth Anything V2 (Metric Depth) 则通过米制深度训练直接预测以米为单位的深度，但预测仍可能存在误差。

因此，尺度信息\textbf{为相对重建补充真实尺度约束，但仅靠一个尺度不能唯一确定整幅图的深度}，各像素的深度仍需预测或测量。本实验分别评估对齐后的几何精度与原始米制精度，以区分结构恢复效果和绝对距离估计能力。
